\documentclass{beamer}
\usepackage{hyperref}
\usepackage[T1]{fontenc}
\usepackage{cite}
\usepackage{relsize}
\usepackage{animate}


\usefonttheme[onlymath]{serif}
%C++ symbol
\newcommand{\Rplus}{\protect\hspace{-.1em}\protect\raisebox{.35ex}{\smaller{\smaller\textbf{+}}}}
\newcommand{\Cpp}{\mbox{C\Rplus\Rplus}\hspace{3pt}}
%Scientific notation
\newcommand{\expnumber}[2]{{#1}\mathrm{e}{#2}}

\usepackage{latexsym,amsmath,xcolor,multicol,booktabs,calligra}
\usepackage{graphicx,tikz,listings,stackengine}
\usetikzlibrary{shapes,snakes,matrix, positioning,arrows.meta}
\setbeamertemplate{bibliography item}[text]

\author{Rafael De Luna Loredo}
\title{Programaci\'on en Python y \Cpp}
\subtitle{Algoritmos de Ordenamiento}
\institute{Facultad de Ciencias, UASLP}
\date{}
\usepackage{NJFU}

% defs
\def\cmd#1{\texttt{\color{red}\footnotesize $\backslash$#1}}
\def\env#1{\texttt{\color{blue}\footnotesize #1}}
\definecolor{deepblue}{rgb}{0,0,0.5}
\definecolor{deepred}{rgb}{0.6,0,0}
\definecolor{deepgreen}{rgb}{0,0.5,0}
\definecolor{halfgray}{gray}{0.55}

\begin{document}
	
	\AtBeginSubsection[]
	{
		\begin{frame}
			\begin{multicols}{3}
				\tableofcontents[currentsection,subsectionstyle=show/shaded/hide,subsubsectionstyle=show/shaded/hide]
			\end{multicols}
		\end{frame}
	}
	\AtBeginSubsubsection[]{
		\begin{frame}
			\begin{multicols}{3}
				\tableofcontents[currentsection,subsectionstyle=show/shaded/hide,subsubsectionstyle=show/shaded/hide]
			\end{multicols}
		\end{frame}
	}
	
	\AtBeginSection[]
	{
		\begin{frame}
			\begin{multicols*}{3}
				\tableofcontents[currentsection,subsectionstyle=show/shaded/hide,subsubsectionstyle=show/shaded/hide]
			\end{multicols*}
		\end{frame}
	}
	
	\begin{frame}
		\titlepage
		\begin{figure}[htpb]
			\begin{center}
				\includegraphics[width=0.2\linewidth]{fig/logo.png}
			\end{center}
		\end{figure}
	\end{frame}
	
	\begin{frame}{Indice}
		\begin{multicols*}{3}
			\tableofcontents[sectionstyle=show,subsectionstyle=show/shaded/hide,subsubsectionstyle=show/shaded/hide]
		\end{multicols*}
	\end{frame}

\section{Introducci\'on}

\begin{frame}{Ordenaci\'on}
	Se refiere al proceso de reagrupar o reorganizar un conjunto de datos u objetos seg\'un un orden o secuencia.
	Veremos que esto nos ser\'a \'util en el siguiente tema.
\end{frame}

\begin{frame}{\textquestiondown Por que ordenar?}
	Imagina que tienes el cl\'oset mas grande del mundo y fuiste invitado a una de las mejores fiestas del mundo, ya tienes en mente que te  vas
	a poner pero toda la ropa est\'a revuelta, hay camisas  donde debería haber pantalones y viceversa o vestidos donde deber\'a haber blusas, todo un caos.
	Pasaste tanto tiempo buscando lo que te ibas a poner que te perdiste de la fiesta.
	Estando ordenando tambi\'en habr\'ia tardado un poco pero sabr\'iamos donde empezar a buscar
\end{frame}

\begin{frame}{Ordenaci\'on y b\'usqueda}
	Aunque ordenar y buscar sean dos cosas muy distintas, est\'an estrechamente relacionadas, ordenar nos ayuda a establecer los l\'imites de la b\'usqueda.
	Si buscamos el n\'umero mas pequeño de un conjunto de diez mil n\'umeros no ordenados, podr\'a ser necesario recorrer todo el conjunto.
	Pero si este estuviera ordenado en forma ascendente o descendente, ni siquiera ser\'ia necesario buscar como tal, sabemos que se encuentra al final o al inicio y
	solo tendr\'iamos que ir a esa posici\'on
\end{frame}

\section{Ordenaci\'on Interna}

\subsection{M\'etodos directos}
\begin{frame}{M\'etodos directos}
	Son m\'etodos sencillos y f\'aciles de comprender pero que son ineficientes cuando los datos a ingresar son demasiados
\end{frame}

\begin{frame}{Ordenaci\'on por intercambio}
	Tambi\'en conocido como m\'etodo burbuja, es el mas f\'acil de implementar pero el mas ineficiente.\\
	Consiste en llevar los elementos mas pequeños hacia la parte izquierda o los mas grandes hacia la parte derecha mediante la
	comparaci\'on de pares adyacentes de elementos e intercambiarlos entre si hasta que todos se encuentren ordenados.
\end{frame}

\subsection{M\'etodos Logar\'itmicos}
\begin{frame}{M\'etodos logar\'itmicos}
	Son m\'etodos mucho mas complejos y dif\'iciles de entender que los directos pero son mas eficientes.
\end{frame}

\section{Referencias}
\subsection{}
\begin{frame}[allowframebreaks]
	
	\frametitle{Referencias}
	\bibliographystyle{acm}
	\bibliography{ref.bib}
\end{frame}

\end{document}
